% Part: first-order-logic
% Chapter: beyond
% Section: modal-logics

\documentclass[../../../include/open-logic-section]{subfiles}

\begin{document}

\olfileid{fol}{byd}{mod}

\olsection{امکانی منطق دیاں قسماں}

ہیٹھلے شرطیہ جملے دی مثال ویکھو:
\begin{quote}
  جے جیریمی اوس کمرے وچ اکیلا اے،
  تاں اوہ نشے وچ، ننگا، تے کرسیاں
  اُتے نچدا پیا اے۔
\end{quote}
ایہ ایداں دا شرطیہ بیان اے جیہڑا
مادی معنی وچ سچ ہو سکدا اے،
پر فیر وی غلط تاثر دیندا اے:
الفاظ ایس گل دا اشارہ کردے نیں
کہ مقدم تے نتیجے وچ اوس توں
مضبوط رشتہ اے جو محض ایہ ہووے
کہ یا مقدم جھوٹا اے یا تالی
سچی اے۔ یعنی جملہ ایہ ظاہر
کردا اے کہ دعویٰ صرف ایس
خاص جہان وچ سچ نہیں (جتھے
ایہ معمولی طور تے سچ ہو
سکدا اے، کیوں جے جیریمی
کمرے وچ اکیلا نہیں)، بلکہ
جے مقدم سچ \emph{ہوندا}
تاں تالی وی \emph{سچ ہوندی}۔
دوجے لفظاں وچ بیان نوں
ایہ سمجھیا جا سکدا اے
کہ دعویٰ صرف ایس جہان وچ
نہیں بلکہ ہر ”ممکن“
جہان وچ سچ اے؛ یا صرف
ایس خاص جہان وچ سچ
ہون دی بجائے
\emph{لازماً} سچ اے۔

امکانی منطق ایس قسم
دی لازمی سچائی دا
مطلب صاف کرن لئی
بنائی گئی اے۔ عام
قضیاتی منطق وچ اک
چوکور عامل شامل
کرکے امکانی قضیاتی
منطق ملدی اے؛ یعنی
جے $!A$ اک
!!a{formula} اے، تاں
$\Box !A$ وی
فارمولا اے۔
بدیہی طور اُتے
$\Box !A$ ایہ
آکھدا اے کہ
$!A$ \emph{لازماً}
سچ اے، یا ہر
ممکن جہان
وچ سچ اے۔
$\Diamond !A$ نوں
عموماً
$\lnot \Box \lnot !A$
دا مختصر نشان
منیا جاندا اے؛ ایس
دا مطلب ایہ
پڑھیا جا سکدا
اے کہ $!A$
\emph{ممکنہ طور
اُتے} سچ اے۔
ظاہر اے
کہ محمولی
منطق وچ
وی امکانی
عامل شامل
کیتے جا
سکدے نیں۔

کرپکے دیاں
!!{structure}s
امکانی منطق
دی معنیات
دین لئی
ورتی جا
سکدیاں نیں؛
اصل وچ
کرپکے نے
ایہ معنیات
پہلاں امکانی
منطق لئی
سوچی سی۔
ایتھے جزوی
ترتیباں تک
محدود رہن
دی بجائے،
عمومی طور
اُتے ”ممکن
جہاناں“ دا
سیٹ $P$
تے جہاناں
وچ دو-جگہی
”رسائی“
دا تعلق
$\Atom{R}{x,y}$
ہوندا اے۔
بدیہی طور
اُتے
$\Atom{R}{p,q}$
آکھدا اے
کہ جہان~$q$
جہان~$p$
نال موافق
اے؛ یعنی
جے اسیں
جہان~$p$
”وچ“ آں،
تاں ایہ
امکان وی
سوچنا پیندا
اے کہ
جہان~$q$
ورگا ہو
سکدا سی۔
کسی جہان وچ چوکور عامل والا بیان اوہدوں سچ ہوندا اے جدوں
اندرلا بیان اوس جہان توں رسائی والے ہر جہان وچ سچ ہووے؛
امکان والے بیان لئی کم از کم اک ایسا جہان کافی اے۔
سو عمومی صورت وچ سارے ممکن جہان نہیں، بلکہ متعلقہ جہان توں
رسائی والے جہان ویکھنے نیں۔ پچھلی وجدانیت پسند معنیات دے برعکس،
ایتھے ایٹمی قضیے دا سچ رسائی دے تعلق نال آگے لازماً برقرار
رہن دی شرط نہیں۔

امکانی منطق
نوں کدی
”مفہومی“
(intensional)
منطق وی
آکھدے نیں،
”امتدادی“
(extensional)
منطق دے
مقابلے وچ۔
کلاسیکی
منطق ورگی
امتدادی
منطق دی
مقصود معنیات
صرف اک
جہان، یعنی
”واقعی“
جہان، دا
حوالہ دیندی
اے؛ جدکہ
مفہومی منطق
دی معنیات
وڈھیری پیچیدہ
وجودیات اُتے
ٹکدی اے۔
لازمیت نوں
!!{structure} دے
روپ وچ پیش
کرن توں علاوہ
امکانی عامل
ہور لسانی
بناوٹاں نوں
!!{structure} دے
روپ وچ پیش
کرن لئی وی
ورتے جا
سکدے نیں:
ورتوں دے
مطابق
$\Box$ تے
$\Diamond$ دا
مطلب نواں
لے لو۔
مثال لئی:
\begin{enumerate}
\item ثابت پذیری
  دی منطق
  وچ $\Box !A$
  نوں ”$!A$
  ثابت ہو
  سکدا اے“
  تے
  $\Diamond !A$
  نوں ”$!A$
  موافق اے“
  پڑھدے نیں۔
  ایہ دوویں پڑھتاں اکّو مقررہ پس منظر دے نظریے دے لحاظ نال نیں:
  ثابت کرن دے قاعدے بدل بدل کے معنی نہیں لینے؛ موافقت دا مطلب
  اوہ نظریہ ایس بیان نال تضاد پیدا نہ کرے۔
\item علم دی
  منطق وچ
  $\Box !A$
  نوں ”میں
  $!A$ جاندا
  آں“ یا
  ”میں $!A$
  منّدا آں“
  پڑھ سکدے
  آں۔
\item زمانی
  منطق وچ
  $\Box !A$
  نوں ”$!A$
  ہمیشہ سچ
  اے“ تے
  $\Diamond !A$
  نوں ”$!A$
  کدی کدی
  سچ اے“
  پڑھیا جا
  سکدا اے۔
\end{enumerate}

اسیں چاہندے
آں کہ منطق
وچ امکانی
عامل بارے
قاعدے تے
مسلّمے شامل
ہون۔ مثال
لئی نظام
$\Log{S4}$
وچ عام
قضیاتی منطق
دے مسلّمے
تے قاعدے،
نال ہیٹھلے
مسلّمے ہوندے
نیں:
% PNB-SEM-034: لازمیت دا قاعدہ بے کھلے مقامی مفروضیاں ثابت نتیجے لئی اے۔
\begin{align*}
& \Box (!A \lif !B) \lif (\Box !A \lif \Box !B)\\
& \Box !A \lif !A \\
& \Box !A \lif \Box \Box !A
\intertext{نال اک قاعدہ وی: جے $!A$ کوئی کھلا مقامی مفروضہ ورتے بغیر
  ثابت شدہ قضیہ اے، تاں $\Box !A$ کڈو۔ ایہ قاعدہ محض فرض کیتے
  ہوئے بیان اُتے نہیں لا سکدے۔
  $\Log{S5}$ ہیٹھلا مسلّمہ شامل کردا اے:}
& \Diamond !A \lif \Box \Diamond !A
\end{align*}
ایہناں مسلّمیاں
دیاں وکھ
صورتاں وکھ
ورتوں لئی
مناسب ہو
سکدیاں نیں؛
مثال لئی
S5 نوں
عموماً منطقی
لازمیت دا
تصور بیان
کرن والا
منیا جاندا
اے۔ چنگی
گل ایہ
اے کہ
کرپکے دیاں
!!{structure}s
وچ رسائی
دے تعلق
نوں فطری
طریقے نال
محدود کرکے
عموماً
ایسی معنیات
مل جاندی
اے جیہدے
لئی !!{derivation}
نظام درست
تے تمام
ہوئے۔
مثال لئی
$\Log{S4}$
انہاں کرپکے
!!{structure}s
دی جماعت
نال ملدا
اے جیہناں
وچ رسائی
دا تعلق
انعکاسی
تے متعدی
اے۔
$\Log{S5}$
انہاں کرپکے
!!{structure}s
دی جماعت
نال ملدا
اے جیہناں
وچ رسائی
دا تعلق
{\em آفاقی}
اے، یعنی
ہر جہان
توں ہر
جہان
تک رسائی
ہوئے؛ ایس
لئی
$\Box !A$
پورا اتردا اے
جے تے
صرف جے
$!A$ ہر
جہان وچ
پورا اترے۔

\end{document}
